\section{Method}
\label{sec:method}

\subsection{Preprocessing}

All articles pass through a shared preprocessing step before any modeling. Each article is split on whitespace, then each token is:

\begin{enumerate}
    \item lowercased,
    \item filtered to retain only tokens matching the pattern \texttt{[\^{}a-zA-Z0-9]} (i.e., alphabetic-initial tokens),
    \item filtered to remove tokens shorter than 2 characters,
    \item optionally filtered against an English stopword list (NLTK \texttt{stopwords.words("english")}).
\end{enumerate}

Stopword removal is applied for genre classification and keyword scoring, where function words carry no topical signal. It is \emph{not} applied for sentiment classification, where negation words (``not,'' ``never'') carry meaningful polarity information.

Unlike many production NLP pipelines, we deliberately avoid stemming or lemmatization. Retaining full word forms keeps the vocabulary human-readable in keyword output and avoids the risk of conflating morphologically similar but semantically distinct terms (e.g., \textit{bank} and \textit{banking} in a business context).

\subsection{Feature Extraction}

\paragraph{Bag-of-Words for classification.}
Both the genre and sentiment classifiers represent documents as bags of words, i.e., unordered multisets of tokens after preprocessing. The preprocessing utility also supports an optional \texttt{binarize} mode, which collapses repeated tokens to presence/absence features; the reported experiments use the default count-based representation (\texttt{BINARIZE=False}).

\paragraph{TF-IDF for keyword ranking.}
For keyword extraction we build an inverted index over all 2,225 articles and compute TF-IDF weights using the standard log-frequency formulation:

\[
\text{tf-idf}(t, d) = (1 + \log_{10} \text{tf}_{t,d}) \times \log_{10}\frac{N}{\text{df}_t}
\]

where $\text{tf}_{t,d}$ is the raw term frequency in document $d$, $N$ is the total number of documents, and $\text{df}_t$ is the number of documents containing term $t$. For articles not seen during index construction (e.g., live RSS articles), we reuse the corpus IDF values and compute TF-IDF on the fly.

\paragraph{POS-tag filtering.}
To produce more readable keyword lists, we additionally filter terms by their NLTK part-of-speech tag, retaining only nouns (NN, NNS, NNP, NNPS) and adjectives (JJ, JJR, JJS). This removes verbs, determiners, and other grammatical function words that might receive high TF-IDF scores in short snippets.

\subsection{Named Entity Recognition}

Named entities are extracted using NLTK's \texttt{ne\_chunk} function \cite{bird2009natural}, which applies a binary Maximum Entropy NE chunker on top of POS-tagged tokens. We extract entities of types PERSON, ORGANIZATION, GPE (geopolitical entity), GSP (geographical/social/political), and FACILITY. Each article's entities are de-duplicated and grouped by type before display.

\subsection{Genre Classifier}

We implement Multinomial Naive Bayes from scratch following the derivation in \citet{mccallum1998comparison}. The model uses bag-of-words counts with Laplace smoothing, estimates class priors from the training set, and skips out-of-vocabulary words at inference time. This keeps the classifier simple, transparent, and effective on the BBC genre labels.

\subsection{Sentiment Analysis}

Sentiment classification uses two independent systems whose outputs are compared:

\begin{itemize}
    \item \textbf{VADER} \cite{hutto2014vader}: a rule-based analyzer using a curated sentiment lexicon with heuristics for capitalization, punctuation, and valence shifters (negation, degree adverbs). It produces a compound score in $[-1, +1]$; we threshold at $\pm 0.05$ to assign positive/negative/neutral labels.

    \item \textbf{NB Sentiment}: the same Multinomial Naive Bayes architecture used for genre classification, trained on BBC articles pseudo-labeled by VADER; articles with compound scores in the neutral band $(-0.05, 0.05)$ are excluded. Training on in-domain text means the model has seen the vocabulary and syntactic patterns of news writing, which we hypothesize should produce more reliable sentiment predictions than a model trained on movie reviews or social media.
\end{itemize}
